\scanpage{21}
% vertical rule at left of the continuation
\noindent\llap{\smash{\raisebox{-0.7\baselineskip}{\rule{0.4pt}{1.6\baselineskip}}}\quad}
isom pour les dim. $p\geq k$ ($k$ entier donné $\leq n$)
il faut et il suffit qu'on ait
\[
H^i(A)=0\qquad\text{pour}\quad k\leq i<n
\]

% vertical rule at left of the corollary
\noindent\llap{\smash{\raisebox{-2.1\baselineskip}{\rule{0.4pt}{2.2\baselineskip}}}\quad}
\uline{Corollaire 1} Pour que ce soit un isom. pour
\deleted{\uline{tout}} $p$, il f. et s. que \inserted{les} $A_{\unsure{m|\mathfrak m}}$ \unsure{soient|soit} \deleted{des}
\deleted{\unsure{modules}} \inserted{anneaux} de Coh. Mac. de dim $n$.

\deleted{\uline{\unsure{Corollaire}}}

\deleted{\uline{Proposition} \unsure{9|2}}\inserted{\uline{Corollaire 2}} Soit $A$ un anneau \inserted{local} de Coh.
Mac. de dim $n$, $\Omega$ un module \deleted{dualisant}\inserted{fond.}
pour $A$. Alors on a
\[
\left\{\begin{aligned}
&\operatorname{Ext}^i(k,\Omega)=0\qquad \text{si}\quad i\ne n\\
&\operatorname{Ext}^n(\unsure{k|h},\Omega)\simeq k
\end{aligned}\right.
\]

\margin{\unsure{\S|$\{$}}
En effet, $\operatorname{Ext}^i(k,\Omega)$ est dual de $H^{n-i}(k)$, qui
est nul si $n-i\ne0$, et qui est $=\underline{k}$ si $n-i=0$

En particulier

\uline{Corollaire 3} $A,\Omega$ comme ci-dessus. Alors $\Omega$
est un module de Coh. Mac. \inserted{de dim $n$} sur $A$, on a
\[
\left\{\begin{aligned}
&H^i(\Omega_{\unsure{A|*}})=0\qquad \text{si}\quad i\ne n\\
&H^n(\Omega_A)\quad\text{est dualisant}
\end{aligned}\right.
\]
De plus, \deleted{$\Omega_A$ est un module fidèle}
\[
\left\{\begin{aligned}
&\operatorname{Ext}^i(\Omega,\Omega)=0\qquad\text{si }i\ne0\\
&\deleted{\operatorname{Ext}}\ \operatorname{Hom}(\Omega,\Omega)\simeq A
\end{aligned}\right.
\]
% archive sheet number: 20

